\documentclass[11pt,a4paper]{article}

% ── Packages ──────────────────────────────────────────────────────────────────
\usepackage[T1]{fontenc}
\usepackage[utf8]{inputenc}
\usepackage{lmodern}
\usepackage[margin=2.5cm]{geometry}
\usepackage{amsmath,amssymb,amsthm}
\usepackage{algorithm}
\usepackage{algpseudocode}
\usepackage{booktabs}
\usepackage{array}
\usepackage{multirow}
\usepackage{xcolor}
\usepackage{hyperref}
\usepackage{microtype}
\usepackage{parskip}
\usepackage{enumitem}
\usepackage{caption}

% ── Colours ───────────────────────────────────────────────────────────────────
\definecolor{coventryblu}{HTML}{1F5C99}
\definecolor{darkgray}{HTML}{444444}
\definecolor{lightblue}{HTML}{EEF4FB}
\definecolor{codegreen}{HTML}{1A6B3C}

\hypersetup{
  colorlinks=true,
  linkcolor=coventryblu,
  citecolor=coventryblu,
  urlcolor=coventryblu
}

% ── Theorem environments ───────────────────────────────────────────────────────
\theoremstyle{plain}
\newtheorem{theorem}{Theorem}[section]
\newtheorem{corollary}[theorem]{Corollary}
\newtheorem{lemma}[theorem]{Lemma}

\theoremstyle{definition}
\newtheorem{definition}[theorem]{Definition}
\newtheorem{example}[theorem]{Example}

\theoremstyle{remark}
\newtheorem{remark}[theorem]{Remark}

% ── Notation macros ───────────────────────────────────────────────────────────
\newcommand{\U}{\mathcal{U}}
\newcommand{\R}{\mathcal{R}}
\newcommand{\E}{\mathcal{E}}
\newcommand{\Sstar}{S^{*}}
\newcommand{\Sprime}{S^{\prime}}
\newcommand{\SCPR}{\textup{SCPR}}
\newcommand{\QUBO}{\textup{QUBO}}
\newcommand{\SCP}{\textup{SCP}}
\newcommand{\ESCP}{\textup{ESCP}}
\newcommand{\abs}[1]{\lvert #1 \rvert}
\newcommand{\set}[1]{\{#1\}}
\newcommand{\minimize}{\operatorname{minimise}}
\newcommand{\subjto}{\operatorname{s.t.}}
\newcommand{\etaR}{\eta_{\mathcal{R}}}
\newcommand{\etaU}{\eta_{\mathcal{U}}}
\newcommand{\etaE}{\eta_{\mathcal{E}}}

% ── Title block ───────────────────────────────────────────────────────────────
\title{%
  \textbf{\color{coventryblu}%
    Encoding the Set Covering Problem with Reasons\\
    as a Quadratic Unconstrained Binary Optimisation Problem:\\
    Equivalence, Beasley Pre-Reduction, and Post-Q Variable Fixing}
}
\author{%
  Abiola Tolulope Babatunde, PhD\\
  \small Centre for Computational Science and Mathematical Modelling\\
  \small Coventry University, United Kingdom\\
  \small \texttt{babatunde.tolulope@outlook.com}
}
\date{June 2025}

% ── Document ──────────────────────────────────────────────────────────────────
\begin{document}

\maketitle

\begin{abstract}
The Set Covering Problem with Reasons (\SCPR{}) generalises the classical Set Covering
Problem by shifting the optimisation objective from the number of covering sets to the
number of underlying \emph{reasons} that activate those sets.
This paper establishes a formal polynomial-time equivalence between \SCPR{} and the
Quadratic Unconstrained Binary Optimisation problem (\QUBO{}), providing a constructive
map in both directions.
Building on Algorithm~2 of~\cite{babatunde2025thesis}---the adapted Beasley reduction
for \SCPR{}---we prove two reduction corollaries:
\emph{pre-Q reduction}, in which Beasley preprocessing shrinks the \QUBO{} variable
count before the matrix is built, reducing qubit requirements on a quantum annealing
processor; and \emph{post-Q variable fixing}, in which Beasley logic is applied directly
to an already-constructed \QUBO{} matrix via forced-variable absorption.
Both approaches preserve optimality under reconstruction $S = \Sprime \cup \Sstar$.
We demonstrate all five stages---\SCPR{} without reduction, \SCPR{} with reduction,
\QUBO{} without reduction, \QUBO{} with pre-Q Beasley, and \QUBO{} with post-Q
Beasley---on three benchmark instances, including the running example from the thesis.
The proof structure mirrors the \ESCP{}$\leftrightarrow$\SCPR{} equivalence
of~\cite{babatunde2025thesis} and positions the \SCPR{} as a native candidate for
quantum annealing via D-Wave's Ocean SDK.
\end{abstract}

\tableofcontents
\newpage

% ──────────────────────────────────────────────────────────────────────────────
\section{Introduction}
\label{sec:intro}
% ──────────────────────────────────────────────────────────────────────────────

The Set Covering Problem (\SCP{}) is a foundational NP-complete combinatorial optimisation
problem with applications in facility location, crew scheduling, network design, and
automated reasoning~\cite{karp1972,beasley1987}.
The \SCPR{} was introduced in~\cite{sadeghimanesh2022smt} and formally analysed
in~\cite{babatunde2025thesis}, extending \SCP{} by associating each covering relationship
with a set of \emph{reasons}: the problem becomes one of minimising the number of reasons
rather than the number of sets.

Separately, Quadratic Unconstrained Binary Optimisation (\QUBO{}) is the native problem
formulation for quantum annealing processors, including D-Wave's Advantage
systems~\cite{dwave2023}.
The \QUBO{} framework has been used to encode a wide range of combinatorial problems,
including scheduling~\cite{pfg2022}, logistics~\cite{portla2022}, and advertising
allocation~\cite{recruit2022}.
Classical set cover is a well-known \QUBO{} encoding target~\cite{glover2018qubo},
but no prior work has established the encoding for the \SCPR{}.

This paper makes the following contributions.

\begin{enumerate}[leftmargin=*, label=\arabic*.]
  \item \textbf{Equivalence theorem} (\S\ref{sec:equivalence}): a constructive
        polynomial-time reduction between \SCPR{} and \QUBO{}, establishing that every
        \SCPR{} instance can be encoded as a \QUBO{} and that every \QUBO{} with
        appropriate structure encodes an \SCPR{}.

  \item \textbf{Pre-Q Beasley reduction} (\S\ref{sec:preq}): application of
        Algorithm~2 of~\cite{babatunde2025thesis} before building the \QUBO{} matrix,
        reducing the variable count and the required penalty coefficient.

  \item \textbf{Post-Q variable fixing} (\S\ref{sec:postq}): application of Beasley
        logic directly to an already-constructed \QUBO{} matrix via forced-variable
        absorption, producing a sparser matrix without rebuilding from scratch.

  \item \textbf{Five-stage computational demonstration} (\S\ref{sec:experiments}):
        all results verified on three benchmark instances.
\end{enumerate}

% ──────────────────────────────────────────────────────────────────────────────
\section{Preliminaries}
\label{sec:prelim}
% ──────────────────────────────────────────────────────────────────────────────

\subsection{The Set Covering Problem with Reasons}

\begin{definition}[\SCPR{}, Definition~1.5.3 of~\cite{babatunde2025thesis}]
\label{def:scpr}
Let $\U = \set{\alpha_1, \ldots, \alpha_n}$ be a universal set,
$\R = \set{\rho_1, \ldots, \rho_r}$ a set of reasons, and
\[
  \E = \set{(A_1, R_1), \ldots, (A_m, R_m)}
\]
a collection of pairs where $A_i \subseteq \U$ are \emph{covering sets} and
$R_i \subseteq \R$ are \emph{reason sets}.
The \SCPR{} is the problem of finding the minimal $S \subseteq \R$ such that
\[
  \bigcup_{\substack{i=1,\ldots,m \\ R_i \subseteq S}} A_i \;=\; \U.
\]
\end{definition}

A pair $(A_i, R_i)$ is \emph{activated} by $S$ if $R_i \subseteq S$.
The \SCPR{} objective is to minimise $\abs{S}$ over all feasible $S \subseteq \R$.

\subsection{Quadratic Unconstrained Binary Optimisation}

\begin{definition}[\QUBO{}]
\label{def:qubo}
Given an upper-triangular matrix $Q \in \mathbb{R}^{n \times n}$,
the \QUBO{} problem is to find $x \in \set{0,1}^n$ minimising the energy
\[
  E(x) \;=\; x^\top Q x \;=\; \sum_{k} Q_{kk}\, x_k \;+\; \sum_{k < j} Q_{kj}\, x_k x_j.
\]
\end{definition}

The diagonal entries $Q_{kk}$ encode linear terms; the off-diagonal entries $Q_{kj}$
encode quadratic interactions.
Constraints are encoded as penalty terms: a violated constraint of coefficient $P$
adds at least $P$ to the energy.

\subsection{The Beasley Reduction for \SCPR{}}

We recall Algorithm~2 of~\cite{babatunde2025thesis}.

\begin{definition}[Beasley Reduction, Algorithm~2 of~\cite{babatunde2025thesis}]
\label{def:beasley}
Given $(\U, \R, \E)$, apply the following two rules iteratively until no further
reduction is possible.
\begin{description}
  \item[Rule~1 (unique coverage).] If $u \in \U$ appears in exactly one pair
        $(A_i, R_i)$, then $R_i \subseteq S$ in every feasible solution.
        Add $R_i$ to $\Sstar$; remove $u$ and all elements of $A_i$ from $\U$;
        remove $R_i$ from $\R$; remove $(A_i, R_i)$ from $\E$.
  \item[Rule~2 (dominance).] If $(A_i, R_i)$ dominates $(A_j, R_j)$---i.e.\
        $A_i \supseteq A_j$ and $R_i \subseteq R_j$---then $(A_j, R_j)$ is not
        needed for an optimal solution. Remove it from $\E$.
\end{description}
The output $(\U', \R', \E', \Sstar)$ satisfies: for any solution $\Sprime$ to
the reduced instance $(\U', \R', \E')$, the set $S = \Sprime \cup \Sstar$ is a
valid solution to the original instance $(\U, \R, \E)$.
\end{definition}

The \emph{reduction efficiency metrics} are defined as
\[
  \etaR = \frac{\abs{\R} - \abs{\R'}}{\abs{\R}},
  \qquad
  \etaU = \frac{\abs{\U} - \abs{\U'}}{\abs{\U}},
  \qquad
  \etaE = \frac{\abs{\E} - \abs{\E'}}{\abs{\E}}.
\]

% ──────────────────────────────────────────────────────────────────────────────
% ──────────────────────────────────────────────────────────────────────────────
\section{\SCPR{} \texorpdfstring{$\leftrightarrow$}{<->} \QUBO{} Equivalence}
\label{sec:equivalence}
% ──────────────────────────────────────────────────────────────────────────────

This section constructs a \QUBO{} whose minimum-energy assignment recovers an
optimal \SCPR{} solution. The construction is the direct \QUBO{} translation of
the optimisation objective already used for the heuristic solvers
in~\cite{babatunde2025thesis} (Chapter~4). In the notation
of~\cite{babatunde2025thesis} that objective maximises
\[
  \phi \;=\; \abs{\R} \;-\; \sum_{k} y_k \;+\; a \sum_{i} \chi_{C_i},
  \qquad a = \abs{\R} + 1,
\]
where $\chi_{C_i} \in \set{0,1}$ is a \emph{saturating} indicator that element
$i$ is covered by \emph{at least one} activated pair. Two properties of that
objective are load-bearing and are preserved here. First, selecting a reason
\emph{increases} cost: in minimisation form the linear weight is $+1$ per reason.
Second, covering an element more than once carries \emph{no} penalty, because the
indicator saturates at $1$. The encoding below realises both as \QUBO{} penalty
terms; departing from either property (a $-1$ objective weight, or a count-based
$P(1-\sum_a a)^2$ coverage term) breaks the equivalence.

\subsection{Variables}

Let $(\U, \R, \E)$ be an \SCPR{} instance with $r = \abs{\R}$. Index the reasons
$\rho_1, \ldots, \rho_r$ with binary variables $y_1, \ldots, y_r$, where
$y_k = 1$ iff $\rho_k \in S$. Partition $\E$ into singleton pairs
$\E_S = \set{(A_i, R_i) : \abs{R_i} = 1}$ and multi-reason pairs
$\E_M = \set{(A_i, R_i) : \abs{R_i} \ge 2}$.

For a singleton pair the \emph{activation variable} is the reason variable $y_k$
itself. For a multi-reason pair $(A_i, R_i) \in \E_M$ the activation is the
conjunction $\bigwedge_{k : \rho_k \in R_i} y_k$. A conjunction of $m \ge 2$
variables is encoded exactly by chaining $m-1$ pairwise \textup{AND} gadgets:
introduce auxiliaries $z_i^{(2)}, \ldots, z_i^{(m)}$ with
$z_i^{(2)} = y_{k_1} \wedge y_{k_2}$ and
$z_i^{(t)} = z_i^{(t-1)} \wedge y_{k_t}$ for $t = 3, \ldots, m$, and take the
final $z_i := z_i^{(m)}$ as the pair's activation variable.

\subsection{Penalties}

Fix a penalty weight $P > \abs{\R}$; the minimum sufficient value is
$P = \abs{\R} + 1$, matching the coefficient $a$ of~\cite{babatunde2025thesis}.

\paragraph{Objective.}
Minimising $\abs{S}$ contributes one linear term per reason,
\begin{equation}
  Q_{kk} \mathrel{+}= 1 \qquad (k = 1, \ldots, r).
  \label{eq:obj}
\end{equation}

\paragraph{Conjunction (\textup{AND}) gadget.}
For each pairwise conjunction $z = a \wedge b$ (where $a, b$ are reason or
intermediate variables) add
\begin{equation}
  P\bigl(ab - 2az - 2bz + 3z\bigr).
  \label{eq:and}
\end{equation}
The polynomial in~\eqref{eq:and} equals $0$ for every assignment consistent with
$z = a \wedge b$ and is strictly positive otherwise. In particular it vanishes
when a strict, possibly nonempty, subset of the conjuncts is selected, so partial
selection of a pair's reasons incurs \emph{no} spurious penalty. (The C1/C2
penalty form $P(z - y_k z)$ and $P(y_k - 2 y_k z + z)$ summed over $k$ does not
have this property: for a partially selected pair it leaves a residual penalty
$P \sum_k y_k$.)

\paragraph{Coverage (at least one).}
For each element $u \in \U$ let $\mathrm{Acts}(u)$ be the activation variables of
the pairs covering $u$, and let $t = \abs{\mathrm{Acts}(u)}$. The covering
requirement is $\sum_{a \in \mathrm{Acts}(u)} a \ge 1$. If $t = 1$ with
activation variable $a$, add
\begin{equation}
  P\,(1 - a),
  \label{eq:cover1}
\end{equation}
which forces $a = 1$. If $t \ge 2$, introduce $L = \lceil \log_2 t \rceil$ binary
slack bits $s_{u,0}, \ldots, s_{u,L-1}$ encoding
$s_u = \sum_{b=0}^{L-1} 2^b s_{u,b} \in \set{0, \ldots, t-1}$, and add
\begin{equation}
  P\Bigl( \textstyle\sum_{a \in \mathrm{Acts}(u)} a \;-\; 1 \;-\; s_u \Bigr)^{2}.
  \label{eq:cover}
\end{equation}
Because $s_u \ge 0$, the penalty~\eqref{eq:cover} attains its minimum (value $0$)
exactly when $\sum_a a \ge 1$, and it is invariant to the number of active pairs
beyond the first: an element covered twice is penalised no more than an element
covered once. This is the \QUBO{} realisation of the saturating indicator
$\chi_{C_u}$ of~\cite{babatunde2025thesis}. Replacing~\eqref{eq:cover} with the
\emph{exact}-cover penalty $P(1 - \sum_a a)^2$ would penalise legitimate
overlapping covers and is the principal source of error in a naive encoding,
since \SCPR{} (Definition~1.5.3 of~\cite{babatunde2025thesis}) is a set cover in
which $\bigcup A_i = \U$ with overlaps permitted.

The full matrix is the sum of~\eqref{eq:obj}, one
copy of~\eqref{eq:and} per pairwise conjunction, and one
of~\eqref{eq:cover1} or~\eqref{eq:cover} per element. Its dimension is
\[
  n \;=\; \abs{\R}
  \;+\; \sum_{(A_i, R_i) \in \E_M} \bigl(\abs{R_i} - 1\bigr)
  \;+\; \sum_{u : t_u \ge 2} \lceil \log_2 t_u \rceil .
\]

\subsection{QUBO to SCPR Construction}

Given $Q$ together with the index partition into reason variables $I_{\R}$,
conjunction auxiliaries $I_A$, and slack bits $I_{\Sigma}$, set
$\R = \set{\rho_k : k \in I_{\R}}$. Each conjunction auxiliary and its incident
couplings identify the reason set $R_i$ of a multi-reason pair; the coverage
penalties~\eqref{eq:cover1}--\eqref{eq:cover} identify, for each element, the
activation variables covering it and hence the covering set $A_i$. A
minimum-energy assignment restricted to $I_{\R}$ yields a minimal feasible
$S \subseteq \R$.

\subsection{Equivalence Theorem}

\begin{theorem}[\SCPR{} $\leftrightarrow$ \QUBO{} Equivalence]
\label{thm:equiv}
Every \SCPR{} instance $(\U, \R, \E)$ can be encoded as a \QUBO{} instance $Q$ of
size
$n = \abs{\R} + \sum_{(A_i,R_i)\in\E_M}(\abs{R_i}-1)
+ \sum_{u : t_u \ge 2}\lceil \log_2 t_u \rceil$
such that a minimum-energy solution of $Q$ recovers an optimal solution of
$(\U, \R, \E)$, and conversely. The encoding is a polynomial-time reduction in
both directions.
\end{theorem}

\begin{proof}
$(\Rightarrow)$ The construction of \S\ref{sec:equivalence} produces $Q$ in time
$O\!\left(\abs{\R} + \sum_i \abs{R_i} + \sum_u t_u\right)$, polynomial in the
instance size.

\textit{Feasibility.}
Let $x^{*}$ be a minimum-energy assignment. Each conjunction
gadget~\eqref{eq:and} contributes $0$ iff its auxiliary equals the conjunction it
represents and at least $P$ otherwise; each coverage
term~\eqref{eq:cover1}--\eqref{eq:cover} contributes $0$ iff the corresponding
element is covered and at least $P$ otherwise. The objective~\eqref{eq:obj} ranges
over $\set{0, \ldots, \abs{\R}}$. Since $P > \abs{\R}$, no objective saving can
offset a single violated penalty, so $x^{*}$ satisfies every conjunction and every
coverage constraint. Its reason bits therefore encode a feasible
$S = \set{\rho_k : x^{*}_k = 1}$ with $\abs{S} = \sum_k x^{*}_k$.

\textit{Optimality.}
Suppose a feasible $S'$ with $\abs{S'} < \abs{S}$ existed. Its canonical encoding
(reason bits from $S'$; each auxiliary set to the induced conjunction; each slack
$s_u$ set to $\sum_{a \in \mathrm{Acts}(u)} a - 1$) satisfies all penalties, so
its energy equals $\abs{S'} < \abs{S} = E(x^{*})$, contradicting minimality of
$x^{*}$.

$(\Leftarrow)$
Given $Q$, the \SCPR{} instance is recovered as in \S\ref{sec:equivalence} in
polynomial time, and a minimum-energy assignment maps to an optimal $S$ by the
same argument. The two directions establish a bijection between feasible solutions
of equal objective value. \hfill $\square$
\end{proof}

\begin{corollary}
The \SCPR{} is NP-complete under \QUBO{} reduction, consistent with
Corollary~2.3.2 of~\cite{babatunde2025thesis}.
\end{corollary}

\begin{remark}[Penalty calibration]
The bound $P > \abs{\R}$ is exactly the condition under which one violated coverage
constraint outweighs deselecting every reason, mirroring Table~4.2
of~\cite{babatunde2025thesis}, where $a = \abs{\R} + 1$ is chosen so that each
uncovered element costs more than any single-reason saving. Because the coverage
penalty~\eqref{eq:cover} encodes an \emph{at-least-one} constraint rather than an
exact count, feasible solutions that overlap incur no penalty, and the energy gap
between competing feasible solutions is governed solely by the
objective~\eqref{eq:obj}, equalling the difference in $\abs{S}$. On quantum
annealing hardware a larger $P$ may be needed to dominate thermal noise; the
pre-reduction of Theorem~\ref{thm:preq} lowers the required ceiling from
$P > \abs{\R}$ to $P > \abs{\R'}$.
\end{remark}

% ──────────────────────────────────────────────────────────────────────────────
\section{Pre-Q Beasley Reduction}
\label{sec:preq}
% ──────────────────────────────────────────────────────────────────────────────

The pre-Q reduction applies Algorithm~2 of~\cite{babatunde2025thesis}
\emph{before} constructing $Q$, yielding a smaller matrix.

\begin{theorem}[Pre-Q Beasley Correctness]
\label{thm:preq}
Let $(\U, \R, \E)$ be an \SCPR{} instance and let
$(\U', \R', \E', \Sstar)$ be the output of the Beasley reduction
(Definition~\ref{def:beasley}).
Let $Q'$ be the \QUBO{} matrix constructed from $(\U', \R', \E')$ via
Theorem~\ref{thm:equiv}.
If $\Sprime \subseteq \R'$ is recovered from a minimum-energy assignment of $Q'$,
then $S = \Sprime \cup \Sstar$ is an optimal solution to the original instance.
\end{theorem}

\begin{proof}
\textit{Feasibility.}
By Rule~1 of Definition~\ref{def:beasley}, every $\rho_k \in \Sstar$ is required
in every feasible solution.
By the reconstruction guarantee of~\cite{babatunde2025thesis}, any $\Sprime$
solving $(\U', \R', \E')$ extends to $S = \Sprime \cup \Sstar$ covering $\U$.
Theorem~\ref{thm:equiv} guarantees that $\Sprime$ solves $(\U', \R', \E')$ iff
the corresponding \QUBO{} assignment is feasible, so $S$ covers $\U$.

\textit{Optimality.}
Suppose for contradiction that a feasible $S^\dagger \subseteq \R$ exists with
$\abs{S^\dagger} < \abs{S}$.
Since $\Sstar \subseteq S^\dagger$ (by necessity of forced reasons), define
$S'^\dagger = S^\dagger \setminus \Sstar \subseteq \R'$.
Then $S'^\dagger$ covers $\U'$ and
$\abs{S'^\dagger} = \abs{S^\dagger} - \abs{\Sstar} < \abs{S} - \abs{\Sstar} = \abs{\Sprime}$,
contradicting optimality of $\Sprime$ on the reduced instance. \hfill $\square$
\end{proof}

\begin{corollary}[Qubit resource reduction]
The pre-Q Beasley reduction reduces the number of reason variables from $\abs{\R}$
to $\abs{\R'} = (1 - \etaR)\abs{\R}$, reducing the number of qubits required for
the \QUBO{} by factor $(1 - \etaR)$ and lowering the required penalty ceiling from
$P > \abs{\R}$ to $P > \abs{\R'}$.
\end{corollary}

\begin{corollary}[Extreme case]
If $\U' = \emptyset$ after the Beasley reduction, then $\Sstar$ is itself an optimal
solution and no \QUBO{} needs to be constructed or submitted to the QPU.
\end{corollary}

% ──────────────────────────────────────────────────────────────────────────────
\section{Post-Q Variable Fixing}
\label{sec:postq}
% ──────────────────────────────────────────────────────────────────────────────

Post-Q variable fixing applies Beasley logic \emph{after} $Q$ has been built,
modifying the matrix directly.
This is useful when $Q$ has already been constructed or received from an external
source and cannot be rebuilt.

\subsection{Variable Fixing in QUBO}

\begin{definition}[Variable fixing]
Given a \QUBO{} matrix $Q \in \mathbb{R}^{n \times n}$ and a set of forced variables
$\set{y_k = 1}$, the \emph{fixed matrix} $Q_{\text{fix}}$ is obtained by:
for each forced index $k$,
\begin{enumerate}[label=(\roman*)]
  \item absorbing the interaction term: $Q_{jj} \mathrel{+}= Q_{kj}$ for all $j \neq k$
        (since $y_k = 1$ reduces $Q_{kj} y_k y_j$ to the linear term $Q_{kj} y_j$);
  \item zeroing out row $k$ and column $k$;
  \item recording the energy offset $Q_{kk}$.
\end{enumerate}
The resulting matrix has $n - 1$ active variables and the same minimum-energy
assignment (restricted to remaining variables) as the original.
\end{definition}

\begin{theorem}[Post-Q Beasley Correctness]
\label{thm:postq}
Let $Q$ be a \QUBO{} encoding of \SCPR{} instance $(\U, \R, \E)$.
Apply Rule~1 of Definition~\ref{def:beasley} to identify forced reasons
$\Sstar \subseteq \R$ and fix the corresponding variables in $Q$ to obtain
$Q_{\text{fix}}$.
If $S_{\text{QUBO}}$ is recovered from a minimum-energy assignment of $Q_{\text{fix}}$,
then $S = S_{\text{QUBO}} \cup \Sstar$ is an optimal solution to $(\U, \R, \E)$.
\end{theorem}

\begin{proof}
By variable fixing, a minimum-energy assignment of $Q_{\text{fix}}$ corresponds
to the minimum-energy assignment of $Q$ restricted to the unfixed variables,
since the fixed variable contributions are absorbed as constant offsets.
The reconstruction $S = S_{\text{QUBO}} \cup \Sstar$ is valid by the same argument
as Theorem~\ref{thm:preq}.
The optimality contradiction argument is identical. \hfill $\square$
\end{proof}

\begin{remark}[Comparison of pre-Q and post-Q reduction]
Pre-Q reduction is preferred when the \QUBO{} has not yet been built, as it
reduces both the number of reason variables and the number of auxiliary variables
(since dominated pairs are also eliminated, removing their auxiliary variables).
Post-Q fixing reduces only the forced reason variables; the auxiliary variables
of eliminated dominated pairs, and of pairs that Rule~1 reduces to a single
reason, remain in $Q_{\text{fix}}$ unless explicitly removed.
The sparsity gain is therefore never larger post-Q than pre-Q, and is strictly
smaller whenever reduction resolves any element entirely, since rebuilding
eliminates the auxiliary and slack variables of resolved pairs that post-Q fixing
leaves in place.
The two approaches always return identical optimal solutions; they differ only in
qubit cost.
In the extreme Rule~1 case of Thesis Example~1.5.2 below, pre-Q resolves the
instance with no \QUBO{} at all, whereas post-Q still constructs and fixes the full
matrix (Table~\ref{tab:results}).
\end{remark}

% ──────────────────────────────────────────────────────────────────────────────
\section{Structural Parallel with ESCP \texorpdfstring{$\leftrightarrow$}{<->} SCPR}
\label{sec:parallel}
% ──────────────────────────────────────────────────────────────────────────────

The proof structure of Theorem~\ref{thm:equiv} is deliberately parallel to
Theorem~2.3.1 of~\cite{babatunde2025thesis}, which establishes the
\ESCP{}$\leftrightarrow$\SCPR{} equivalence.
Both proofs follow the same three-step pattern.

\begin{enumerate}[leftmargin=*, label=\arabic*.]
  \item \textbf{Constructive map}: a polynomial-time transformation of one problem
        representation into the other (\SCPR{} $\to$ \ESCP{}; \SCPR{} $\to$ \QUBO{}).

  \item \textbf{Forward feasibility}: any solution to the source representation maps
        to a feasible solution of the target with equal objective value.

  \item \textbf{Backward contradiction}: a strictly better target solution would imply
        a strictly better source solution, contradicting optimality.
\end{enumerate}

The key structural difference is that the \ESCP{}$\leftrightarrow$\SCPR{} equivalence
operates at the combinatorial problem level (rewriting set families), whereas the
\SCPR{}$\leftrightarrow$\QUBO{} equivalence operates at the computational substrate
level (encoding into a real-valued matrix).
Both are polynomial-time reductions; both preserve NP-completeness.

% ──────────────────────────────────────────────────────────────────────────────
\section{The Five-Stage Pipeline Algorithm}
\label{sec:algorithm}
% ──────────────────────────────────────────────────────────────────────────────

Algorithm~\ref{alg:pipeline} summarises the complete five-stage pipeline.

\begin{algorithm}[H]
\caption{SCPR Five-Stage \QUBO{} Pipeline}
\label{alg:pipeline}
\begin{algorithmic}[1]
\Require \SCPR{} instance $(\U, \R, \E)$; penalty coefficient $P > \abs{\R}$
\Ensure Optimal reason set $S \subseteq \R$

\Statex \textbf{Stage~1:} Solve \SCPR{} directly (brute-force reference)
\State $S_{\text{ref}} \gets$ brute-force enumeration of subsets of $\R$ \Comment{reference solution}

\Statex \textbf{Stage~2:} Apply Beasley reduction to \SCPR{}
\State $(\U', \R', \E', \Sstar) \gets \mathrm{Beasley}(\U, \R, \E)$
\If{$\U' = \emptyset$} \Return $\Sstar$ \EndIf
\State $\Sprime \gets$ brute-force on $(\U', \R', \E')$
\State $S \gets \Sprime \cup \Sstar$ \Comment{reconstruction; verify $S$ covers $\U$}

\Statex \textbf{Stage~3:} Build and solve \QUBO{} without reduction
\State $Q \gets \mathrm{BuildQUBO}(\U, \R, \E, P)$
\State $S_3 \gets \mathrm{Solve}(Q)$

\Statex \textbf{Stage~4:} Pre-Q Beasley reduction then \QUBO{}
\State $(\U', \R', \E', \Sstar) \gets \mathrm{Beasley}(\U, \R, \E)$
\If{$\U' = \emptyset$} \Return $\Sstar$ \EndIf
\State $Q' \gets \mathrm{BuildQUBO}(\U', \R', \E', P')$ where $P' > \abs{\R'}$
\State $\Sprime \gets \mathrm{Solve}(Q')$
\State $S_4 \gets \Sprime \cup \Sstar$ \Comment{Theorem~\ref{thm:preq}}

\Statex \textbf{Stage~5:} Build full \QUBO{} then post-Q variable fixing
\State $Q \gets \mathrm{BuildQUBO}(\U, \R, \E, P)$
\State $(\Sstar, Q_{\text{fix}}) \gets \mathrm{PostQFix}(Q, \U, \R, \E)$ \Comment{Definition~5.1}
\State $S_{\text{QUBO}} \gets \mathrm{Solve}(Q_{\text{fix}})$
\State $S_5 \gets S_{\text{QUBO}} \cup \Sstar$ \Comment{Theorem~\ref{thm:postq}}
\end{algorithmic}
\end{algorithm}

% ──────────────────────────────────────────────────────────────────────────────
\section{Experiments}
\label{sec:experiments}
% ──────────────────────────────────────────────────────────────────────────────

We demonstrate all five stages on three benchmark instances.
All experiments ran in Python~3.12.3 with NumPy~2.4.4, SciPy~1.17.1 and
\texttt{dimod}~0.12.23.
Each \QUBO{} stage finds its minimum-energy assignment by exhaustive
enumeration, and brute-force enumeration of subsets of $\R$ supplies the
reference optimum.
On the three benchmarks, \texttt{dimod}'s \texttt{ExactSolver} and
\texttt{SimulatedAnnealingSampler} (200 reads) reach the same minimum energy as
the enumeration.
A randomised suite of $4{,}000$ feasible instances, with up to 5 elements,
5 reasons and 7 pairs (seed~1), checks Stage~2 on every instance and the
\QUBO{} Stages~3 to~5 on the $3{,}356$ instances whose \QUBO{} has at most
20 variables; the other 644 exceed the enumeration cap.
No stage disagreed with the brute-force optimum.
Source code and the full harness output are available at
\url{https://github.com/ToluPeazy/scpr-qubo}.

\subsection{Instance Descriptions}

\textbf{Small\_4x4.}
$\abs{\U} = 4$, $\abs{\R} = 4$, $\abs{\E} = 4$.
No Beasley reduction triggers; serves as a baseline showing the pipeline is
transparent when preprocessing has no effect.

\textbf{Reduction\_Demo.}
$\abs{\U} = 5$, $\abs{\R} = 4$, $\abs{\E} = 4$.
Element~4 appears in exactly one pair; Rule~1 forces $r_2$ into $\Sstar$.
Designed to isolate the effect of a single Rule~1 application.

\textbf{Thesis\_1.5.2.}
The running example from~\cite{babatunde2025thesis} (16~lights, 8~batteries).
Rule~1 applies twice: elements~4 and~9 are each covered by exactly one pair.
Beasley solves the instance entirely ($\U' = \emptyset$).

\subsection{Results}

Table~\ref{tab:results} summarises the five stages across all three instances.

\begin{table}[H]
\centering
\caption{Five-stage results. All solutions are feasible and optimal.}
\label{tab:results}
\small
\setlength{\tabcolsep}{4pt}
\begin{tabular}{llccccc}
\toprule
Instance & Stage & $\abs{S}$ & Covers~$\U$ & QUBO vars & $\etaR$ & Notes \\
\midrule
\multirow{5}{*}{Small\_4x4}
& 1: Brute-force        & 3 & \checkmark & -- & 0\% & Reference \\
& 2: SCPR+reduction     & 3 & \checkmark & -- & 0\% & $\Sstar = \emptyset$ \\
& 3: QUBO no-red        & 3 & \checkmark & 11 & 0\% & $P = 5$ \\
& 4: QUBO pre-Q         & 3 & \checkmark & 11 & 0\% & No reduction \\
& 5: QUBO post-Q        & 3 & \checkmark & 11 & 0\% & No variables fixed \\
\midrule
\multirow{5}{*}{Reduction\_Demo}
& 1: Brute-force        & 3 & \checkmark & -- & --  & Reference \\
& 2: SCPR+reduction     & 3 & \checkmark & -- & 25\% & $\Sstar = \set{r_2}$ \\
& 3: QUBO no-red        & 3 & \checkmark & 10 & 0\% & $P = 5$ \\
& 4: QUBO pre-Q         & 3 & \checkmark & 8  & 25\% & $P = 4$; $-$2 vars \\
& 5: QUBO post-Q        & 3 & \checkmark & 9  & 25\% & $P = 5$; $-$1 var \\
\midrule
\multirow{5}{*}{Thesis\_1.5.2}
& 1: Brute-force        & 6 & \checkmark & -- & --  & Reference \\
& 2: SCPR+reduction     & 6 & \checkmark & -- & 75\% & $\Sstar = \set{0,1,2,3,4,6}$ \\
& 3: QUBO no-red        & 6 & \checkmark & 21 & 0\% & $P = 9$ \\
& 4: QUBO pre-Q         & 6 & \checkmark & -- & 75\% & Solved; no QUBO needed \\
& 5: QUBO post-Q        & 6 & \checkmark & 15 & 75\% & $-$6 vars fixed \\
\bottomrule
\end{tabular}
\end{table}

The variable counts follow the construction of Section~\ref{sec:equivalence}:
$\abs{\R}$ reason variables, $\sum_i (\abs{\R_i} - 1)$ auxiliary variables from the
chained pairwise \textup{AND} gadgets, and $\sum_u \lceil \log_2 t_u \rceil$ coverage
slack variables, contributing one block per element coverable by $t_u \geq 2$ pairs.
The slack and chained-auxiliary variables are the cost of a faithful set-cover
encoding that admits legitimate overlaps; a naive exact-cover penalty would use fewer
qubits but does not preserve the \SCPR{} optimum (Section~\ref{sec:equivalence}).

\subsection{Discussion}

The results confirm three claims.

First, Theorem~\ref{thm:equiv} holds empirically: the Stage~3 \QUBO{}
minimum-energy solutions match the brute-force reference in Stage~1 on all three
benchmarks and on all $3{,}356$ random instances small enough to enumerate, with
zero discrepancies.

Second, pre-Q reduction (Stage~4) saves one more variable than post-Q fixing
(Stage~5) on Reduction\_Demo: 8 variables against 9, from 10 without reduction.
Rule~1 forces $r_2$, which turns the pair with reason set $\set{r_1, r_2}$ into a
single-reason pair, so the \QUBO{} rebuilt on the reduced instance needs no
auxiliary for it.
Post-Q fixing removes only $y_{r_2}$; the auxiliary $z = y_{r_1} \wedge y_{r_2}$
stays in $Q_{\text{fix}}$, where it now tracks $y_{r_1}$.
Rule~2 removes no pairs on this instance, so the saving comes from Rule~1 alone.
Rebuilding after reduction removes auxiliaries that variable fixing leaves in
place, as the remark in \S\ref{sec:postq} notes.

Third, the extreme-case corollary applies to Thesis~1.5.2: both Stage~2 and
Stage~4 identify $\U' = \emptyset$ after two Rule~1 applications, returning
$\Sstar$ immediately without constructing or solving any \QUBO{}.
This is the most computationally efficient outcome: the full 75\% reason reduction
eliminates the need for QPU submission entirely.

% ──────────────────────────────────────────────────────────────────────────────
\section{Future Work}
\label{sec:future}
% ──────────────────────────────────────────────────────────────────────────────

\begin{enumerate}[leftmargin=*, label=\arabic*.]
  \item \textbf{QPU validation.} Submit reduced \QUBO{} instances to D-Wave
        Advantage via Ocean SDK and measure solution quality as a function of
        chain-break rate and annealing time.

  \item \textbf{Penalty calibration.} Derive instance-specific bounds on $P$
        as a function of $\abs{\R'}$, energy gap, and QPU coupling precision.

  \item \textbf{Lagrangian post-Q reduction.} Extend post-Q fixing beyond
        Rule~1 to include LP relaxation-based variable elimination
        (Beasley's original Lagrangian bound), applied directly to $Q$.

  \item \textbf{Benchmark extension.} Apply the five-stage pipeline to the
        full \SCPR{} benchmark datasets of~\cite{babatunde2025thesis}
        (Chapters~5--7), reporting QUBO-specific metrics alongside the existing
        $\etaR$, $\etaU$, $\etaE$ tables.

  \item \textbf{CipherRescue application.} Apply the pipeline to FDE failure
        diagnosis \SCPR{} instances from the CipherRescue framework,
        targeting USENIX Security~/ IEEE S\&P.
\end{enumerate}

% ──────────────────────────────────────────────────────────────────────────────
\section{Conclusion}
\label{sec:conclusion}
% ──────────────────────────────────────────────────────────────────────────────

We have established a formal polynomial-time equivalence between the \SCPR{} and
\QUBO{}, extending the equivalence programme of~\cite{babatunde2025thesis} to the
quantum computing substrate.
The constructive proof provides an explicit encoding that correctly maps \SCPR{}
solutions to minimum-energy \QUBO{} assignments and vice versa.
Two reduction corollaries---pre-Q Beasley and post-Q variable fixing---provide
practical mechanisms for reducing QPU resource requirements without sacrificing
optimality, with the pre-Q approach being strictly more efficient when the instance
has reducible auxiliary variables.
The five-stage demonstration confirms all theoretical claims computationally on
three benchmark instances, including the Thesis~1.5.2 example where the 75\%
reason reduction eliminates the need for QPU submission entirely.

The \SCPR{}'s cross-domain applicability---demonstrated in~\cite{babatunde2025thesis}
across SMT solving, cancer genomics, and cryptographic failure diagnosis---suggests
that these \QUBO{} encoding results have broad practical reach as quantum annealing
hardware continues to scale.

% ──────────────────────────────────────────────────────────────────────────────
\bibliographystyle{plain}
\begin{thebibliography}{99}

\bibitem{babatunde2025thesis}
A.~T.~Babatunde.
\newblock Introducing and Analysing the Set Covering Problem with Reasons.
\newblock PhD thesis, Coventry University, 2025.

\bibitem{sadeghimanesh2022smt}
A.~Sadeghimanesh and M.~England.
\newblock An SMT solver for non-linear real arithmetic inside Maple.
\newblock Extended abstract, 2022.

\bibitem{karp1972}
R.~M.~Karp.
\newblock Reducibility among combinatorial problems.
\newblock In \textit{Complexity of Computer Computations}, pages 85--103.
  Plenum Press, 1972.

\bibitem{beasley1987}
J.~E.~Beasley.
\newblock An algorithm for set covering problem.
\newblock \textit{European Journal of Operational Research}, 31(1):85--93, 1987.

\bibitem{glover2018qubo}
F.~Glover, G.~Kochenberger, and Y.~Du.
\newblock Quantum Bridge Analytics I: a tutorial on formulating and using QUBO models.
\newblock \textit{4OR}, 17:335--371, 2019.

\bibitem{dwave2023}
D-Wave Systems.
\newblock D-Wave Ocean SDK documentation.
\newblock \url{https://docs.dwavequantum.com}, 2023.

\bibitem{pfg2022}
D-Wave Systems.
\newblock Quantum application in scheduling: Pattison Food Group.
\newblock Technical report, D-Wave Quantum Inc., 2022.

\bibitem{portla2022}
D-Wave Systems.
\newblock Quantum application in logistics: SavantX and Port of Los Angeles.
\newblock Technical report, D-Wave Quantum Inc., 2022.

\bibitem{recruit2022}
D-Wave Systems.
\newblock Quantum application in advertising: Recruit Communications.
\newblock Technical report, D-Wave Quantum Inc., 2022.

\end{thebibliography}

\end{document}

